\scanpage{109}
% archive foliation: 108
\textcircled{2} Cas où on dispose seulement d'une représentation de $F$ sous forme sympathique \uline{localement} sur $X$ pour la top (Zar). Dans ce cas,
\deleted{on ne peut plus utiliser \unsure{comme} \unsure{puisque} $X$ séparé, \unsure{on doit en faire en localisant} sur $X$ \unsure{un sens} (\unsure{ét}), et utiliser la suite spectrale de Leray d'un recouvrement. On trouve \unsure{un énon. d'effectivité}, poids entre $0$ et $2n+2j$}
on a intérêt à travailler \uline{d'abord} avec la cohomologie à support propre sur $X$, qui permet de découper $X$ en morceaux, et d'obtenir les $=$ résultats que dessus. On déduit l'analogue en cohomologie à supports quelconques par dualité.

\deleted{Si $F$ est sympathique}

\uline{Remarque} [En fait, on voit par ce raisonnement que, en ce qui concerne la coh. à support propre, les résultats sont obtenus en supposant seulement que $F$ soit ``joli par morceaux'' sur $X$; \inserted{la lissité de $X$ par Poincaré}. Le fait $F$ \unsure{const.|const. et} tordu \inserted{sur $X$ lisse} est alors utilisé \deleted{\unsure{dans le cas lisse}} pour pouvoir procéder par dualité.

\textcircled{C} \uline{Cas de coefficients effectifs de (poids $j$, niveau $j$)} $j$ \uline{\unsure{quelc.}};
$X$ quelconque.

1) Cas de la cohomologie $H^i_!$ : déjà traité dans la remarque précédente : résultat le $=$ que dans \textcircled{B}.1).

2) Cas de la cohomologie $H^i(X,F)$. Je dis qu'il est \uline{effectif}, de poids \deleted{\unsure{entre $0$ et $2i$}} \inserted{$\le\operatorname{Inf}(2i,2n)=2\operatorname{Inf}(i,n)$}, de niveau $\le\operatorname{Inf}(\unsure{2i|i},n)$, où $n=\dim X$. \deleted{\unsure{Cas propre pour}} \ldots

\margin{?}
